\section{Endpoints and analysis plan}
\label{sec:analysis}

\paragraph{Primary endpoints.}
The product-performance endpoint is C4 mean one-shot execution accuracy across
three sealed repetitions, estimated from 303 question-repetition attempts with
question-clustered uncertainty. C4 repetition-one accuracy is reported separately
as the intuitive single-invocation view. There is no majority vote, best-of-three,
or selective-retry score. The comparative endpoint is the paired C4 minus C1
execution-accuracy difference on the same trials, paired within question and
repetition and clustered by question. It is a system-level contrast between two
reasonably developed systems.

\paragraph{Exploratory contrasts.}
C2 minus C1, C3 minus C2, and C4 minus C3 are prespecified exploratory
mechanistic evidence, reported with effect sizes, transition counts, and
uncertainty. Modest rung differences will not be promoted on significance. Where
C3 and C4 lack model or system parity, C4 minus C3 is explicitly system-level.
Repetition-one McNemar analyses may appear as sensitivity checks; naive McNemar
will not be applied to 303 correlated observations. With 101 questions, small
paired deltas will be imprecise, and effect sizes and failure transitions take
precedence over p-values.

\paragraph{Repetitions and reliability.}
Each condition runs three independent repetitions on every test question, for
1,212 frozen outputs generated before any scoring. Trial order is a committed
deterministic block-interleaved permutation over question, condition, and
repetition, with repetitions of one question separated, and each trial records
timestamps and observable model identifiers. The reliability analysis reports
per-repetition and mean accuracy; per-question 3/3, 2/3, 1/3, and 0/3
correctness; correctness-flip rate separately from result-set changes; refusal and
error rate and its stability; latency and cost distributions; and
question-clustered bootstrap intervals that resample questions as clusters
retaining all condition and repetition outcomes. Correctness stability and answer
stability are reported separately: the same wrong answer three times is stable and
incorrect, while different SQL returning the same correct result is
outcome-reliable. SQL-text variation is diagnostic only.

\paragraph{Co-outcomes.}
Alongside accuracy we report confidently-wrong rate, refusal and error rate,
tokens per correct answer, median and interquartile range for tokens and latency,
tool calls and database queries per attempt and per correct answer, retry and
validation distributions, telemetry coverage, and terminal failure vectors by
condition. Operational numerators include every valid attempt; comparative views
use matched question-condition-repetition populations. Cost comparisons require
equivalent output accounting, since a cheaper condition that fails to answer is
not more efficient. Each major intervention reports whether it increased
correctness, reduced confident errors, converted wrong answers into refusals, or
only shifted the failure mode, so a flat accuracy result can still be
product-relevant.

\paragraph{Subgroups.}
Prespecified descriptive subgroups are database, public \texttt{high\_level}, and
public \texttt{conditions} values with adequate cell size. Rare cells and
per-database test estimates are descriptive and will not be mined for
significance. Relative improvement is reported only where the baseline
denominator is nonzero.

\paragraph{Failure taxonomy.}
Failure categories are built from training evidence rather than forced into a
prewritten list, and each records count, denominator, representative non-private
examples, detection evidence, and a primary ownership hypothesis (model, semantic
model, harness, product, environment, evaluator, or benchmark) that changes only
with evidence. For a failed development question whose offline annotation
references HKB knowledge, we classify the earliest supported failure point on a
mechanism ladder: knowledge absent from the semantic model; knowledge present but
the dependency graph wrong; knowledge correct but not retrieved; retrieved but
misinterpreted; representation correct but compilation failed; compiled query
correct but validation or harness behavior changed the outcome; or model
reasoning failed despite a correct available representation. Hidden knowledge
identifiers stay in offline diagnosis, and the published taxonomy records
aggregate classifications only.

Prioritization uses the qualitative product of prevalence, plausible fixability,
and generality, with a model or reviewer supplying the semantic judgment and its
evidence; application code validates record structure and enforces custody policy
rather than embedding keyword heuristics for failure meaning.

\paragraph{Retry and rerun policy.}
The evaluated system's normal production retry policy is part of the frozen
system. Exhausting it, returning no answer, generating invalid SQL, or timing out
inside the agent or harness is an incorrect system outcome. A trial may be rerun
only when failure is demonstrably outside the evaluated system: sealed-evaluator
process failure, benchmark database unavailability, or transport failure before
the evaluated service accepted the request. A rerun keeps the same trial
identity and appends a reason with both attempt records. Rate limits, 5xx
responses, and timeouts are classified in advance by which system owns the retry
boundary, and are never reclassified after scoring. No trial is ever rerun
because its answer was wrong.
